\documentclass[12pt]{article}
\usepackage[utf8]{inputenc}
\usepackage[T1]{fontenc}
\usepackage{hyperref}
\title{Khao Language Resources for K-12 Curriculum Development}
\author{Compiled by Claude Sonnet 5.5 (claude-sonnet-5-5) for Andr\'as Kornai}
\date{}
\begin{document}
\maketitle

\section{Introduction}
Khao (ISO 639-3 xao, Glottocode khao1244) is an Austroasiatic language of northwestern Vietnam. The ``Bookkeeping'' family label in the prompt appears to be a data error, since the reference description says Austro-Asiatic. Glottolog treats Khao as an older name for part of the K\'hang (Kh\'ang) people, with the same language as K\'hang [kjm]. K\'hang is placed in Khmuic or Palaungic depending on the author. Wikipedia gives 14,000 speakers (2009 census) for K\'hang, and a Lingweb entry estimates about 10,000 speakers in S\a'on La and \DH i\d{e}n Bi\^en provinces. Khao should not be confused with the Kho Mu, a different group. The searches found no Khao-specific educational materials, so the resources below are reference and background only.

\section{Primary Educational Resources}
None were found. No textbook, primer or script for Khao or K\'hang turned up, and nothing indicates Khao-medium schooling.

\section{Linguistic and Reference Materials}
\begin{itemize}
\item Glottolog entry for Khao [xao]: \url{https://glottolog.org/resource/languoid/iso/xao}
\item K\'hang language (Wikipedia): \url{https://en.wikipedia.org/wiki/Kh%C3%A1ng_language}
\item Lingweb (MPI-EVA) Khao numeral system entry: \url{https://lingweb.eva.mpg.de/channumerals/Khao.htm}
\end{itemize}

\section{Religious and Community Resources}
Joshua Project profile of the Khao in Vietnam: \url{https://joshuaproject.net/people_groups/12653}. It is a people-group profile and not a language resource. No scripture or community organisation was identified.

\section{Digital and Audio Resources}
None were found. A UNESCO-related item titled ``Documenting and Preserving the Khang Language in Vietnam'' was cited in a reference list, but no URL or content was seen.

\section{Curriculum Development Context}
Little exists. Vietnamese is the language of schooling, and there is no evidence of a Khao writing system or of mother-tongue instruction. A curriculum would need new documentation work with the community, starting from the K\'hang linguistic literature and Vietnamese ethnic-minority education policy, which these searches did not cover.

\end{document}